\begin{textsample}{POS Dim 1 – rj – Score 32.00 – rj/rj\_ef\_ciencias\_humanas\_his\_3.txt}  \label{ex:f1_pos_251}
História no Ensino Fundamental Considerando \textbf{que} o ensino de História no Ensino Fundamental está ligado diretamente a questões como “ o processo de constituição do \textbf{sujeito} ” e seus \textbf{desdobramentos} no \textbf{posicionamento} individual e coletivo no tempo e no espaço, em toda a sua complexidade, conforme o \textbf{que} nos é apresentado de forma \textbf{detalhada} no texto introdutório da BNCC, História - Anos Iniciais, vamos perceber o sentido da gradação da complexidade crescente do \textbf{que} está proposto na própria organização do documento em si e a preocupação com a necessidade de estarmos pautados em \textbf{posicionamentos} éticos e de respeito à diversidade como condição \textbf{inerente} à \textbf{humanidade}.
A cada ano de escolaridade, professores e estudantes vão aprofundando mutuamente a construção de seus conhecimentos da história, considerando \textbf{que} a dinâmica do espaço-tempo presente e suas implicações vão trazendo, a cada etapa do Ensino Fundamental, novas possibilidades de exploração da realidade histórica a nossa volta.
Assim, também quando formos colocar em prática o nosso planejamento curricular, seja no âmbito da nossa cidade, nossa escola ou mesmo em cada sala de aula, devemos cuidar para \textbf{que} cada passo seja dado levando em conta a diversidade de \textbf{sujeitos}, visões, culturas e pensamentos presentes na realidade \textbf{imediata} com \textbf{que} estamos lidando.
É com este olhar \textbf{que} trabalharemos as unidades temáticas, os objetos de conhecimento e habilidades \textbf{aqui} propostas, sabendo \textbf{que} na efetivação de \textbf{um} currículo, outros aspectos terão de ser considerados, tais como as estratégias, os procedimentos \textbf{metodológicos}, as ferramentas de avaliação processual do aprendizado, entre outras questões \textbf{que} somente serão contempladas na discussão coletiva e na reflexão docente no seu cotidiano escolar específico.
Destacamos as unidades temáticas apontadas para cada ano de escolaridade na etapa dos anos iniciais do Ensino Fundamental : - Mundo pessoal : meu lugar no mundo ; - Mundo pessoal : eu, meu grupo social e meu tempo ; - A comunidade e seus registros ; - As formas de registrar as experiências da comunidade ; - O trabalho e a sustentabilidade na comunidade ; - As pessoas e os grupos \textbf{que} compõem a cidade e o município ; - O lugar em \textbf{que} \textbf{vive} ; - A noção de espaço público e privado ; - Transformações e permanências nas trajetórias dos grupos humanos ; - Circulação de pessoas, produtos e culturas ; - As questões históricas relativas às migrações ; - Povos e culturas : meu lugar no mundo e meu grupo social ; - Registros da história : linguagens e culturas.
Se nos anos iniciais do ensino fundamental, processo de ensino/aprendizagem de história deve se \textbf{basear} no reconhecimento de noções como “ Eu ”, “ Nós ” e “ Os outros ”, nos anos finais do ensino fundamental, esse processo, \textbf{que} é qualitativamente diferente, \textbf{baseia} -se em procedimentos como a identificação e organização de eventos importantes na história do ocidente, ressaltando -se a realidade brasileira, a realidade do estado do Rio de Janeiro, as regionalidades e as especificidades locais, com a subsequente ordenação desses conhecimentos de forma cronológica e dentro de seus respectivos espaços geográficos.
Além disso, os estudantes devem proceder à análise e produção de documentos, elaborando \textbf{uma} crítica desses registros, mostrando a capacidade de utilizar diversas linguagens para tanto.
Devem também conhecer/reconhecer diversas perspectivas de análise dos fenômenos históricos, tendo acesso e podendo refletir sobre essas diferentes perspectivas de forma consciente, crítica e autônoma.
Parte das habilidades presentes nesta proposta já está presentes no documento da Base Nacional Comum Curricular e outra parte incorpora as sugestões dadas para inclusão no DOC-RJ elaborado de maneira alinhada à BNCC.
Estas unidades, e suas referentes habilidades podem ser trabalhadas no período de \textbf{um} ano letivo, mas deve -se sempre \textbf{atentar} para a necessidade da progressão e da continuidade do conhecimento.
Assim, qualquer habilidade ou unidade temática pode ser revisitada em anos posteriores, de acordo com as necessidades impostas pela realidade de cada rede, de cada escola, e de cada sala de aula.
Assim, a partir do 6º ano do ensino fundamental, as habilidades requeridas têm como ponto focal a reflexão sobre o significado da história, suas formas de registro, seus \textbf{métodos} de análise e a construção de seu saber, assim como a reflexão sobre as primeiras formas de organização social da \textbf{humanidade}, avançando sobre a Antiguidade clássica e o período medieval.
Em nossa proposta, as hipóteses referentes para a ocupação do território brasileiro, e especificamente de nosso estado, devem ganhar certo destaque.
Para isso amparamo -nos nas mais recentes pesquisas realizadas por arqueólogos como Walter Neves e Niéde Guidon, \textbf{que} oferecem perspectivas diferentes para a chegada dos seres humanos, não só ao Brasil, mas ao continente americano como \textbf{um} todo, contradizendo \textbf{teorias} eurocêntricas sobre o mesmo assunto.
Em consonância com as habilidades, inicialmente se apresentam as seguintes unidades temáticas dos anos finais : História - tempo, espaço, metodologia e formas de registro ; A História, o tempo, e a construção do \textbf{homem} ; Política, organização do poder e relações sociais na antiguidade ; O mundo do trabalho na antiguidade ; e A cultura, o pensar e as pessoas no mundo antigo, tentando abarcar todos os eixos temáticos pertinentes a nossa \textbf{disciplina}, metodologia, poder e ideias, economia e sociedade, e cultura e mentalidades.
A proposta é enfatizar as relações de poder presentes na Europa medieval, a ruptura da lógica feudal a partir do \textbf{século} XIV, e o processo de Expansão Marítima e construção do mundo colonial, enfatizando se tratar de \textbf{um} primeiro momento de “ globalização ” ( ou mundialização do capital se preferirmos ), \textbf{que} se deu entre os \textbf{séculos} XV e XVIII, e deixou profundas \textbf{marcas} em todos esses lugares, sobretudo na América e na África.
No caso de currículo das escolas das redes municipais e estaduais do Rio de Janeiro, é \textbf{interessante} enfatizar os diferentes povos \textbf{que} habitavam o litoral fluminense, destacando -se o papel de tribos como os Tupinambás, Aimorés, Goytacazes (... ), \textbf{que} tiveram papel importante nos primeiros anos da colonização portuguesa no Brasil, mas \textbf{que}, sobretudo, eram povos \textbf{que} já habitavam a região antes da chegada dos portugueses, e \textbf{que} marcaram o território com sua forma de organização social, a utilização dos recursos naturais e sua cultura.
Ressaltar o papel do Brasil e, em especial, o papel do estado do Rio de Janeiro, dentro dessa dinâmica é fundamental ao longo dos anos finais, em especial no primeiro ano deste segmento.
Entretanto, é \textbf{preciso} romper com a visão simplista de \textbf{que} o Brasil era apenas \textbf{um} complemento da sua metrópole.
Demonstrar o Brasil como \textbf{um} lugar \textbf{que} possuía \textbf{uma} dinâmica própria, tanto em sua sociedade como em sua economia, e, sobretudo em sua cultura, e \textbf{que} muitas vezes essa dinâmica se contrapunha a dinâmica proposta pela metrópole, pode servir para desconstruirmos essa visão simplista.
Para tanto, analisar a sociedade colonial, principalmente nos \textbf{séculos} XVI e XVII, como algo muito mais complexo \textbf{que} a dicotomia senhor/escravos, é \textbf{um} desafio para essas séries.
Exemplos como a “ Revolta da cachaça ” em 1660, podem ser ilustrativos para \textbf{frisar} esses momentos em \textbf{que} os interesses dos colonos se confrontavam com os interesses reinos.
Outro ponto a ser ressaltado, é a vida dos escravizados, tanto africanos como indígenas, e de como, a partir de \textbf{uma} situação de opressão, esses \textbf{conseguiram} manter, de diversas formas, suas tradições, e ao mesmo tempo, \textbf{conseguiram} também fazer com \textbf{que} essas se incorporassem às tradições do povo brasileiro em geral.
Na sequência, apresentam -se as unidades temáticas Poder, riqueza e “ espiritualidade ” no mundo medieval ; O mundo \textbf{moderno} e a conexão entre sociedades africanas, americanas e europeias ; Humanismos, Renascimentos e o Novo Mundo ; A organização do poder e as dinâmicas do mundo colonial americano ; e Lógicas comerciais e mercantis da modernidade.
A seguir, ressalta -se a formação do mundo \textbf{contemporâneo}, a partir dos processos impulsionados pela burguesia europeia desde os \textbf{séculos} XVII e XVIII, mas sobretudo, é fundamental perceber como esses processos chegaram \textbf{aqui} no continente americano, e como estes processos adquiriram toda \textbf{uma} dinâmica diferente ao chegarem às Américas.
Nesse caso, é \textbf{interessante}, no caso do Brasil, entender como se deu o processo de “ emancipação política ” do país, como consequência da \textbf{crise} do sistema colonial português, mas também como consequência das revoltas anticoloniais do fim do XVIII e início do XIX, \textbf{que} assumiram por \textbf{aqui}, \textbf{um} caráter muito mais nativista e antilusitano do \textbf{que} propriamente nacionalista, na acepção \textbf{que} o termo tem \textbf{hoje}. \textbf{Focar} na construção do Estado brasileiro, já no XIX, com seus limites e avanços, e com, sobretudo, o entendimento \textbf{que} grande parte da população ( de origem negra, indígena e pobre ), não tinha acesso a direitos básicos e ao exercício da cidadania, e \textbf{que}, quando se revoltavam, eram reprimidos com extrema dureza e crueldade ; é essencial para entendermos as origens da nossa construção como país, e nesse sentido, o Rio de Janeiro, como sede do império, vai \textbf{desempenhar} \textbf{um} papel decisivo.
Dividimos as unidades temáticas referentes a este ano, em quatro : O mundo \textbf{contemporâneo} - o Antigo Regime em \textbf{crise} ; Os processos de independência nas Américas, O Brasil no \textbf{século} XIX, Configurações do mundo no \textbf{século} XIX.
Porém destrinchamos essas unidades em habilidades bem extensas, \textbf{que} abarcam \textbf{uma} gama bem abrangente de fenômenos.
Finalmente, as temáticas a serem desenvolvidas no último ano, giram em torno das questões \textbf{que} perpassaram o \textbf{século} XX ( muitas delas se arrastando pelo \textbf{século} XXI ), tanto no Brasil como no mundo.
Questões como democracia x ditadura, totalitarismo, nacionalismos, extremismos, revoluções, utopias ( e o suposto “ fim ” destas ), movimentos de libertação, perseguições políticas, conflitos étnicos, genocídios, guerras, terrorismo, avanço tecnológico, mudanças políticas e econômicas, mundo \textbf{globalizado}, etc....
No Brasil, destaca -se o combate as ditaduras do \textbf{século} XX ( Vargas e Militar ), e os distintos projetos políticos em jogo naquele período, e a longa luta pela construção de \textbf{um} Estado democrático no país, o \textbf{que} engloba diversas demandas \textbf{que} ainda não estão concretizadas por completo, além dos assuntos relacionados especificamente com a história do Rio de Janeiro, como o “ bota abaixo ” de Pereira Passos, a política de remoção de “ favelas ” do governo Lacerda, as origens da criminalidade no estado, etc.
As unidades temáticas desse ano são quatro : O nascimento da República no Brasil e os processos históricos até a metade do \textbf{século} XX ; Totalitarismos e conflitos mundiais ; Modernização, ditadura civil-militar e redemocratização - o Brasil após 1946 ; e A história recente.
Enfim, o fio condutor desta proposta curricular, é a discussão de temáticas pertinentes ao nosso tempo, priorizando os elementos de história local mais ligados a vida de nossos alunos, para \textbf{que} \textbf{possamos} atender as demandas e aos conhecimentos necessários para \textbf{um} \textbf{posicionamento} cidadão perante as questões \textbf{que} nos são apresentadas no mundo de \textbf{hoje}.
Desta forma se apresenta o Documento Curricular do estado do Rio de Janeiro, no componente curricular história, com o propósito de trazer para a pauta docente o debate frequente sobre a importância da formação de todos os \textbf{sujeitos} envolvidos na educação de crianças e jovens, na perspectiva de \textbf{que} estes estejam, coletiva e individualmente se preparando para “ enfrentar os desafios do mundo \textbf{contemporâneo} ”.
As Competências Específicas de História para o Ensino Fundamental Assim como a BNCC, a proposta de termos \textbf{um} Documento Curricular para o estado do Rio de Janeiro leva em conta \textbf{que} professores e estudantes venham a ser sujeitos do próprio processo de \textbf{ensino-aprendizagem}, sem deixar de considerar \textbf{que} todos esses últimos tenham a garantia do desenvolvimento das competências estabelecidas para a Área de Ciências Humanas em articulação, como é proposto, com as específicas deste componente curricular para o ensino fundamental, abaixo descritas : 1.
Compreender acontecimentos históricos, relações de poder e processos e mecanismos de transformação e manutenção das estruturas sociais, políticas econômicas e culturais ao longo do tempo e em diferentes espaços para analisar, posicionar -se e intervir no mundo \textbf{contemporâneo}. 2.
Compreender a historicidade no tempo e no espaço, relacionando acontecimentos e processos de transformação e manutenção das estruturas sociais, políticas, econômicas e culturais, bem como problematizar os significados das lógicas de organização cronológica. 3.
Elaborar questionamentos, hipóteses, argumentos e proposições em relação a documentos, interpretações e contextos históricos específicos, recorrendo a diferentes linguagens e mídias, exercitando a empatia, o diálogo, a resolução de conflitos, a cooperação e o respeito. 4.
Identificar interpretações \textbf{que} expressem visões de diferentes \textbf{sujeitos}, culturas e povos com relação a \textbf{um} mesmo contexto histórico, e posicionar -se criticamente com base em princípios éticos, democráticos, inclusivos, sustentáveis e solidários. 5.
Analisar e compreender o movimento de populações e mercadorias no tempo e no espaço e seus significados históricos, levando em conta o respeito e a solidariedade com as diferentes populações. 6.
Compreender e problematizar os conceitos e procedimentos \textbf{norteadores} da produção historiográfica. 7.
Produzir, avaliar e utilizar tecnologias digitais de informação e comunicação de modo crítico, ético e responsável, compreendendo seus significados para os diferentes grupos ou estratos sociais.

% matched lemmas: aqui, atentar, basear, conseguir, contemporâneo, crise, desdobramento, desempenhar, detalhar, disciplina, ensino-aprendizagem, focar, frisar, globalizado, hoje, homem, humanidade, imediato, inerente, interessante, marca, metodológico, moderno, método, norteador, posicionamento, possar, preciso, que, sujeito, século, teoria, um, viver
\end{textsample}
